\subsection{THE EULER-MACLAURIN SUMMATION FORMULA}

In the formula (35) of VI, p.~280, let us replace x by 0 and h by x; since \(\mathrm{B}_{m}(0)=b_{m}\) it follows from the relations (24) of VI, p.~276, that for each integer p\textgreater0 one can write

\[
\begin{array}{l} f (x) = \int_ {x} ^ {x + 1} f (t) d t - \frac {1}{2} \big (f (x + 1) - f (x) \big) \\ \qquad + \sum_ {k = 1} ^ {p} \frac {b _ {2 k}}{(2 k) !} \left(f ^ {(2 k - 1)} (x + 1) - f ^ {(2 k - 1)} (x)\right) + \mathrm{R} _ {p} (x) \end{array}\tag{37}
\]

with

\[
\mathrm{R} _ {p} (x) = - \frac {1}{(2 p + 1) !} \int_ {0} ^ {1} \mathrm{B} _ {2 p + 1} (t) f ^ {(2 p + 1)} (x + 1 - t) d t.\tag{38}
\]

In this formula let us successively replace \(x\) by \(x + 1, x + 2, \ldots, x + n\) , and combine the formulae obtained, one by one; we obtain

\[
\left\{ \begin{array}{l} f (x) + f (x + 1) + \dots + f (x + n) \\ = \int_ {x} ^ {x + n + 1} f (t) d t - \frac {1}{2} \bigl (f (x + n + 1) - f (x) \bigr) \\ \quad + \sum_ {k = 1} ^ {p} \frac {b _ {2 k}}{(2 k) !} \bigl (f ^ {(2 k - 1)} (x + n + 1) - f ^ {(2 k - 1)} (x) \bigr) + \mathrm{T} _ {p} (x, n) \end{array} \right.\tag{39}
\]

with

\[
\mathrm{T} _ {p} (x, n) = - \frac {1}{(2 p + 1) !} \int_ {0} ^ {1} \mathrm{B} _ {2 p + 1} (t) \left(\sum_ {k = 0} ^ {n} f ^ {(2 p + 1)} (x + k + 1 - t)\right) d t.\tag{40}
\]

The remainder \(T_{p}(x,n)\) in this formula can be rewritten in the following way: denote by \(\overline{\mathbf{B}}_{2p+1}(t)\) the periodic function with period 1 which is equal to \(\mathbf{B}_{2p+1}(t)\) on the interval {[}0, 1{[}. Then

\[
\int_ {0} ^ {1} \mathrm{B} _ {2 p + 1} (t) f ^ {(2 p + 1)} (x + k + 1 - t) d t = \int_ {k} ^ {k + 1} \overline {{\mathrm{B}}} _ {2 p + 1} (1 - s) f ^ {(2 p + 1)} (x + s) d s
\]

and consequently

\[
\mathrm{T} _ {p} (x, n) = - \frac {1}{(2 p + 1) !} \int_ {0} ^ {n + 1} \overline {{\mathrm{B}}} _ {2 p + 1} (1 - s) f ^ {(2 p + 1)} (x + s) d s.\tag{41}
\]

The formula (39) is called the Euler-Maclaurin summation formula; it is applicable to every complex function having a continuous \((2p+1)^{th}\) derivative on an interval \([x_{0},+\infty[\) , for every \(x\geqslant x_{0}\) . We shall see (VI, p.~288) how to estimate the remainder \(T_{p}(x,n)\) in this formula.

